\section{The arithmetic model remains projective over the modeled affine algebra}

\begin{lemma}[A relative projective smooth arithmetic family]
\label{lem:kn-relative-arithmetic-original-geometric-model}
\lean{EnvelopeProperty.Vanishing.RelativeArithmeticGeometricModel.exists_original_geometric_model}
\uses{lem:kn-relative-arithmetic-original-projective-stage,lem:kn-relative-arithmetic-ambient-base-change}
\leanok
In the situation of the original-projective-stage lemma, there are
an integer subalgebra $C\subset k$ of finite type, a finitely
presented $C$-algebra $D$ with $k\otimes_C D\simeq R$, and an
actual cartesian original-field square $X=Y\times_C k$ such that
$Y\to\operatorname{Spec}C$ is smooth of relative dimension $n$
and separated.  The same $Y$ has a projective morphism
$Y\to\operatorname{Spec}D$ whose composite to
$\operatorname{Spec}C$ is its original coefficient structure.
There is a very ample line $T$ for this relative morphism whose
whole pullback to the original $X$ is
$j^*\mathcal O_{\mathbb P_R^d}(1)$.
\end{lemma}
\begin{proof}
First fix a smooth closed arithmetic stage; the preceding stage producer
supplies this stage from the original data.  Use its smooth closed model
$Y$ in a coefficient extension of $\mathbb P_B^d$, and let $C$ be its actual finite
coefficient subalgebra.  Put $D=C\otimes_A B$.  Finite
presentation is stable under this base change, and tensor
associativity, followed by the original $k\otimes_A B\simeq R$,
gives $k\otimes_C D\simeq R$.  The ambient base-change lemma
gives a relative projective map from
$\mathbb P_B^d\times_A C$ to $\operatorname{Spec}D$, with very
ample pulled-back $\mathcal O(1)$.  Identify the actual ambient
stage with this cartesian product by its two original projections.
Compose the produced closed embedding of $Y$ with that
identification.  Restriction of the very ample line is very ample
on $Y$; in particular the resulting map to
$\operatorname{Spec}D$ is projective.  Its coefficient triangle
is the original coefficient triangle because each cartesian
identification preserves both projections.

The ambient map to $\operatorname{Spec}C$ is separated: it is
the base change of the composite of projective space over $B$
with the affine map $\operatorname{Spec}B\to\operatorname{Spec}A$.
Its closed subscheme $Y$ is therefore separated over $C$.
Finally, the original recovery equation for the closed embedding
identifies the original pullback of $T$ with the pullback of
$\mathcal O(1)$ along $j$ followed by the original coefficient map
$\mathbb P_R^d\to\mathbb P_B^d$.  The whole projective twist
base-change isomorphism identifies this with the stated original
$j^*\mathcal O(1)$.

The stage coefficient algebra is first regarded as an $A$-subalgebra
of $k$.  Restricting its scalars to $\mathbb Z$ gives the stated
integer subalgebra with the same underlying ring.  Finite type is
transitive over $\mathbb Z\to A$, and this identification preserves
the spectrum, the tensor algebra, all four maps of the original
cartesian square, and the entire recovered line isomorphism.
\end{proof}
